\documentclass[runningheads]{llncs}

% ---------------------------------------------------------------
% Include basic ECCV package

% TODO REVIEW: Insert your submission number below by replacing '*****'
% TODO FINAL: Comment out the following line for the camera-ready version
\usepackage[review,year=2026,ID=*****]{eccv}
% TODO FINAL: Un-comment the following line for the camera-ready version
%\usepackage{eccv}

% ---------------------------------------------------------------
% Other packages

\usepackage{eccvabbrv}
\usepackage{graphicx}
\usepackage{booktabs}
\usepackage{amsmath,amssymb}
\usepackage{algorithm}
\usepackage{algorithmic}
\usepackage{multirow}
\usepackage{subcaption}
\usepackage{xcolor}
\usepackage{tikz}
\usetikzlibrary{arrows.meta,positioning,fit,backgrounds,calc,shapes.geometric}

% The "axessiblity" package can be found at: https://ctan.org/pkg/axessibility?lang=en
\usepackage[accsupp]{axessibility}

% ---------------------------------------------------------------
% Hyperref package
% TODO FINAL: Comment out the following line for the camera-ready version
\usepackage[pagebackref,breaklinks,colorlinks,citecolor=eccvblue]{hyperref}
% TODO FINAL: Un-comment the following line for the camera-ready version
%\usepackage{hyperref}
\usepackage[capitalize]{cleveref}

% Support for ORCID icon
\usepackage{orcidlink}

% Note: \eg, \ie, \etal are already defined by eccvabbrv

\begin{document}

% ---------------------------------------------------------------
\title{UrbanAlign: Post-hoc Semantic Calibration of VLM-Synthesised Urban Perception Data}

\titlerunning{UrbanAlign}

\author{Anonymous ECCV 2026 Submission}
\authorrunning{Anonymous}
\institute{Anonymous Institution}

\maketitle

% ===============================================================
% ABSTRACT
% ===============================================================
\begin{abstract}
Large vision-language models (VLMs) can describe urban scenes fluently but produce poor pairwise-comparison labels when used as zero-shot annotators.
We argue that the bottleneck is not the VLM's perceptual ability but its \emph{uncalibrated mapping} from rich internal representations to discrete judgement labels.
We propose \textbf{UrbanAlign}, a post-hoc calibration framework that routes VLM outputs through interpretable semantic dimensions and a locally-adaptive calibration layer.
Semantic dimensions discovered by the VLM form a compact information bottleneck; multi-agent deliberation extracts robust continuous scores; and locally-weighted ridge regression (LWRR) on a hybrid visual--semantic manifold calibrates these scores against human ratings.
An end-to-end optimization loop automatically selects the best-performing dimension set across temperature-scheduled trials.
On Place Pulse~2.0 across all six perception categories (safety, beautiful, lively, wealthy, boring, depressing), UrbanAlign achieves 61.3\% accuracy ($\kappa{=}0.22$) on average---a +17.9\,pp gain over zero-shot VLM (43.4\%) and +7.0\,pp over the uncalibrated Mode~4 baseline (54.3\%)---while providing dimension-level interpretability for every prediction.

\keywords{Urban Perception \and Post-hoc Calibration \and Vision-Language Model \and Semantic Dimensions \and Local Manifold Regression}
\end{abstract}


% ===============================================================
% 1  INTRODUCTION
% ===============================================================
\section{Introduction}
\label{sec:intro}

Urban perception---the subjective assessment of environmental qualities such as safety, wealth, and aesthetics---is a key input to urban planning, real-estate valuation, and social-equity analysis~\cite{lynch1960image,jacobs1961death}.
The Place Pulse project~\cite{salesses2013collaborative} collected 1.17M pairwise comparisons across 110K street-view images, enabling data-driven perception research.
Subsequent work trained deep networks to predict these labels~\cite{dubey2016deep,zhang2018measuring}, but at enormous annotation cost: covering 10K locations at 10 raters each costs ${\sim}\$167$K.

Large VLMs (\eg, GPT-4o~\cite{openai2024gpt4o}) offer a seemingly cheap alternative.
However, benchmarks reveal a striking gap: Mushkani~\etal~\cite{mushkani2025llm} report only 31\% macro-average accuracy across 30 urban-perception dimensions; Beneduce~\etal~\cite{beneduce2025exploring} find $\text{F1}{=}59.2$\% on Place Pulse.
In our own experiments, a zero-shot GPT-4o baseline achieves 43.4\% accuracy ($\kappa{=}0.10$) on average across six categories---better than random but far from the human inter-rater consistency of $\kappa{\approx}0.6$.

\smallskip\noindent\textbf{Why do VLMs fail as annotators yet succeed as scene describers?}
We hypothesise that the problem is not perceptual capability but \emph{calibration}: VLMs can extract rich, relevant features from urban scenes (fa\c{c}ade quality, vegetation condition, pavement integrity), yet their direct-label outputs are poorly aligned with the specific decision boundaries that human raters implicitly apply.
This is analogous to the observation in Concept Bottleneck Models (CBMs)~\cite{koh2020concept} that prediction accuracy improves dramatically when raw features are routed through \emph{interpretable concept scores} and then a \emph{calibrated linear layer}, rather than used end-to-end.
Yuksekgonul~\etal~\cite{yuksekgonul2023posthoc} extended this idea post-hoc: a concept layer can be retrofitted onto any trained backbone without retraining.

UrbanAlign (\cref{fig:framework}) operationalises this principle for VLM-synthesised data.
Rather than treating VLM outputs as final labels, we treat them as \emph{uncalibrated concept scores}---a semantic bottleneck---and learn a lightweight, locally-adaptive calibration function against a small human-labelled reference set.
The key insight is that VLMs are strong \emph{feature extractors} (producing dimension scores that capture genuine perceptual information) but poor \emph{decision makers} (mapping those features to discrete left/right/equal labels).
Post-hoc calibration corrects the mapping while preserving the rich semantic content.

\subsection{Contributions}

\begin{enumerate}
\item \textbf{Post-hoc semantic calibration framework.}
    We propose UrbanAlign, a three-stage pipeline that repositions VLMs from black-box annotators to structured semantic decoders, with a lightweight LWRR calibration stage that aligns VLM concept scores to human perception.
    An end-to-end dimension optimization loop automatically selects the highest-performing dimension set via temperature-scheduled trials with elite seeding, achieving convergence probability $\geq 1-(1-p^{*})^{T}$.

\item \textbf{Interpretable semantic dimensions.}
    Semantic dimensions extracted by the VLM (Fa\c{c}ade Quality, Vegetation Maintenance, Pavement Integrity, \etc) achieve up to 64.4\% individual discriminative power on \emph{wealthy} perception.

\item \textbf{Multi-agent deliberation for robust feature extraction.}
    An Observer--Debater--Judge chain, combined with pairwise image presentation, yields a +13.3\,pp accuracy gain over single-shot scoring on \emph{wealthy} perception.
    A $2{\times}2$ factorial ablation reveals that the two factors interact synergistically: multi-agent reasoning is effective only with pairwise context.

\item \textbf{Local manifold calibration with interpretability audit.}
    LWRR on a hybrid CLIP--semantic differential manifold improves Mode~4 accuracy by +7.3\,pp.
    The local $R^{2}$ provides a per-sample audit: when $R^{2}$ is high, abstract perception is linearly explainable by mid-level dimensions.
\end{enumerate}


% ---------------------------------------------------------------
% Figure 1: Method comparison (three-pipeline conceptual diagram)
% ---------------------------------------------------------------
\begin{figure}[t]
\centering
\definecolor{pipeagray}{HTML}{6C7A89}
\definecolor{pipebgray}{HTML}{8E9EAB}
\definecolor{pipeours}{HTML}{1A5276}
\definecolor{oursaccent}{HTML}{16A085}
\definecolor{badred}{HTML}{C0392B}
\definecolor{goodgreen}{HTML}{27AE60}
\resizebox{\textwidth}{!}{%
\begin{tikzpicture}[
    >=Stealth,
    node distance=0.4cm and 0.5cm,
    pbox/.style={rectangle, draw=#1, fill=#1!8, rounded corners=3pt,
                 minimum height=1.0cm, align=center, font=\small},
    inputbox/.style={rectangle, draw=gray!70, fill=gray!5, rounded corners=3pt,
                 minimum height=1.2cm, text width=2.0cm, align=center, font=\small},
    verdict/.style={font=\small\bfseries, #1},
    arr/.style={->, thick, #1},
    stagelabel/.style={font=\footnotesize\itshape, text=gray!70},
]

% ---- Shared input on the left ----
\node[inputbox, text width=2.2cm, minimum height=3.8cm] (input) {%
  \includegraphics[width=1.6cm,height=1.2cm,draft]{streetview.jpg}\\[2pt]
  {\footnotesize Urban Street}\\[-2pt]{\footnotesize View Image}};

% ========== Pipeline A: End-to-End / Zero-Shot VLM ==========
\node[pbox=pipeagray, right=1.0cm of input, yshift=1.4cm, text width=2.4cm] (pa1)
    {End-to-End\\Zero-Shot VLM};
\node[pbox=pipeagray, right=0.5cm of pa1, text width=2.2cm] (pa2)
    {{\large ?}\\[-2pt]\footnotesize Black Box\\[-2pt]\footnotesize Processing};
\node[pbox=pipeagray, right=0.5cm of pa2, text width=2.8cm] (pa3)
    {\footnotesize ``Image A looks\\[-2pt]\footnotesize more wealthy''\\[-1pt]
     \tiny (No explanation)};
\node[verdict=badred, right=0.3cm of pa3] (pav) {\texttimes};
\node[right=0.05cm of pav, font=\footnotesize, text=badred, text width=2.2cm] {Low\\[-2pt]Interpretability};

\draw[arr=pipeagray] (input.east |- pa1) -- (pa1);
\draw[arr=pipeagray] (pa1) -- (pa2);
\draw[arr=pipeagray] (pa2) -- (pa3);

% ========== Pipeline B: Low-Level Objective ==========
\node[pbox=pipebgray, right=1.0cm of input, text width=2.4cm] (pb1)
    {Low-Level\\Feature Extraction};
\node[pbox=pipebgray, right=0.5cm of pb1, text width=2.2cm] (pb2)
    {\footnotesize Pixel Counting\\[-2pt]\footnotesize Segmentation\\[-2pt]
     \tiny (DeepLab / NDVI)};
\node[pbox=pipebgray, right=0.5cm of pb2, text width=2.8cm] (pb3)
    {\footnotesize Greenery: 30\%\\[-2pt]\footnotesize Building: 45\%\\[-1pt]
     \tiny Acc.\ 59.1\%};
\node[verdict=badred, right=0.3cm of pb3] (pbv) {\texttimes};
\node[right=0.05cm of pbv, font=\footnotesize, text=badred, text width=2.2cm] {Limited\\[-2pt]Semantic Depth};

\draw[arr=pipebgray] (input.east |- pb1) -- (pb1);
\draw[arr=pipebgray] (pb1) -- (pb2);
\draw[arr=pipebgray] (pb2) -- (pb3);

% ========== Pipeline C: UrbanAlign (Ours) — highlighted ==========
\node[pbox=pipeours, right=1.0cm of input, yshift=-1.4cm, text width=2.4cm,
      line width=1.2pt, fill=pipeours!12] (pc1)
    {\textbf{UrbanAlign}\\[-1pt]\footnotesize\textit{(Ours)}};
\node[pbox=pipeours, right=0.5cm of pc1, text width=2.2cm,
      line width=1.2pt, fill=pipeours!12] (pc2)
    {\footnotesize Mid-Level\\[-2pt]\footnotesize Semantic\\[-2pt]\footnotesize Decoding\\[-1pt]
     \tiny\color{oursaccent} VLM + Multi-Agent};

% Dimension score bars
\node[pbox=pipeours, right=0.5cm of pc2, text width=2.8cm,
      line width=1.2pt, fill=pipeours!12] (pc3) {%
  \footnotesize Fa\c{c}ade:\enspace
    {\color{oursaccent}\rule{0.65cm}{5pt}}~8.2\\[-1pt]
  \footnotesize Pavement:
    {\color{oursaccent}\rule{0.50cm}{5pt}}~6.5\\[-1pt]
  \footnotesize Cleanliness:
    {\color{oursaccent}\rule{0.33cm}{5pt}}~4.2\\[-1pt]
  \tiny Acc.\ \textbf{72.3\%}
};
\node[verdict=goodgreen, right=0.3cm of pc3] (pcv) {\checkmark};
\node[right=0.05cm of pcv, font=\footnotesize, text=goodgreen, text width=2.4cm] {Accurate \&\\[-2pt]Interpretable};

\draw[arr=pipeours, line width=1.2pt] (input.east |- pc1) -- (pc1);
\draw[arr=pipeours, line width=1.2pt] (pc1) -- (pc2);
\draw[arr=pipeours, line width=1.2pt] (pc2) -- (pc3);

% Row labels
\node[stagelabel, above=0.02cm of pa1.north west, anchor=south west] {(a)};
\node[stagelabel, above=0.02cm of pb1.north west, anchor=south west] {(b)};
\node[stagelabel, above=0.02cm of pc1.north west, anchor=south west] {(c)};

\end{tikzpicture}%
}
\caption{Conceptual comparison of three approaches to urban perception.
\textbf{(a)}~End-to-end / zero-shot VLM methods map directly from pixels to abstract judgments with no interpretable intermediates.
\textbf{(b)}~Low-level objective methods (segmentation-based regression) extract hand-crafted features but lack semantic depth (59.1\% avg.\ accuracy).
\textbf{(c)}~UrbanAlign routes predictions through mid-level semantic dimensions discovered by the VLM (\eg, Fa\c{c}ade Quality, Pavement Integrity), achieving 72.3\% accuracy with actionable, per-dimension interpretability.}
\label{fig:comparison}
\end{figure}

% ===============================================================
% 2  RELATED WORK
% ===============================================================
\section{Related Work}
\label{sec:related}

\noindent\textbf{Urban perception and crowdsourcing.}
Salesses~\etal~\cite{salesses2013collaborative} created Place Pulse, collecting pairwise comparisons of urban scenes across six dimensions.
Dubey~\etal~\cite{dubey2016deep} trained Siamese networks on these labels.
Zhang~\etal~\cite{zhang2018measuring} used multi-task learning; Naik~\etal~\cite{naik2014streetscore} predicted continuous safety scores.
Abascal~\etal~\cite{abascal2024correlation} extended the paradigm to remote-sensing imagery.
All these methods map directly from pixels to abstract percepts, bypassing interpretable mid-level representations.

\smallskip\noindent\textbf{VLMs for urban understanding.}
Mushkani~\etal~\cite{mushkani2025llm} benchmarked eight VLMs on 30 urban-perception dimensions (best: 31\% macro accuracy).
Beneduce~\etal~\cite{beneduce2025exploring} explored persona-based prompting on Place Pulse ($\text{F1}{=}59.2$\%).
Both position VLMs as end-to-end classifiers.
We reposition VLMs as \emph{structured semantic decoders} whose outputs require post-hoc calibration.

\smallskip\noindent\textbf{Concept Bottleneck Models.}
Koh~\etal~\cite{koh2020concept} introduced CBMs: predictions are routed through interpretable concept scores, enabling human-auditable reasoning.
Yuksekgonul~\etal~\cite{yuksekgonul2023posthoc} proposed post-hoc CBMs that retrofit concept layers onto any backbone.
Oikarinen~\etal~\cite{oikarinen2023labelfree} showed concepts can be discovered from language models without manual labels.
Yang~\etal~\cite{yang2023language} further explored language-in-a-bottle concept extraction.
Our semantic dimensions serve as an analogous interpretable bottleneck, but operate \emph{externally} on VLM-synthesised data rather than as an internal network layer.

\smallskip\noindent\textbf{Multi-agent LLM reasoning.}
Wang~\etal~\cite{wang2023selfconsistency} showed self-consistency (sampling diverse reasoning paths) improves chain-of-thought reliability.
Du~\etal~\cite{du2024improving} and Liang~\etal~\cite{liang2024encouraging} extended this to multi-agent debate.
We instantiate this as an Observer--Debater--Judge pipeline for visual perception.

\smallskip\noindent\textbf{Locally-weighted regression.}
Cleveland~\cite{cleveland1979robust} introduced LOESS for robust local fitting; Atkeson~\etal~\cite{atkeson1997locally} generalised it.
Ruppert and Wand~\cite{ruppert1994multivariate} analysed the multivariate case.
We adapt locally-weighted ridge regression to a hybrid visual--semantic differential space for perception calibration.

\smallskip\noindent\textbf{Few-shot data synthesis.}
Tang~\etal~\cite{tang2025urbanalign} explored LLM-based annotation synthesis with distribution alignment.
Ye~\etal~\cite{ye2022zerogen} and Meng~\etal~\cite{meng2023tuning} generated training data from language models.
UrbanAlign differs by operating on visual pairwise comparisons and introducing post-hoc calibration of VLM concept scores.


% ===============================================================
% 3  METHOD
% ===============================================================
\section{Method}
\label{sec:method}

% ---------------------------------------------------------------
% Figure 2: Framework overview (TikZ)
% ---------------------------------------------------------------
\begin{figure*}[t]
\centering
\definecolor{vlmblue}{HTML}{1A5276}
\definecolor{clipteal}{HTML}{16A085}
\definecolor{feedred}{HTML}{A93226}
\resizebox{\textwidth}{!}{%
\begin{tikzpicture}[
    >=Stealth,
    node distance=0.6cm and 0.8cm,
    block/.style={rectangle, draw=#1, fill=#1!8, rounded corners=3pt,
                  minimum height=0.9cm, text width=2.8cm, align=center,
                  font=\small},
    smallblock/.style={rectangle, draw=#1, fill=#1!8, rounded corners=2pt,
                  minimum height=0.7cm, text width=2.0cm, align=center,
                  font=\footnotesize},
    agent/.style={rectangle, draw=vlmblue, fill=vlmblue!12, rounded corners=2pt,
                  minimum height=0.65cm, text width=1.6cm, align=center,
                  font=\footnotesize},
    data/.style={rectangle, draw=gray, fill=gray!8, rounded corners=2pt,
                  minimum height=0.6cm, align=center, font=\footnotesize},
    stage/.style={rectangle, draw=#1, fill=#1!15, rounded corners=4pt,
                  minimum height=1.0cm, font=\small\bfseries, align=center},
    arr/.style={->, thick, #1},
    darr/.style={->, thick, dashed, #1},
]

% ========== STAGE 1 ==========
\node[stage=vlmblue, text width=3.2cm] (s1label) {\color{vlmblue}Stage 1\\[-1pt]\footnotesize Semantic Dimension\\[-2pt]\footnotesize Extraction};

\node[data, right=0.6cm of s1label, text width=2.2cm] (consensus) {Consensus\\Samples\\{\tiny ($\mu$-high / $\mu$-low)}};
\node[block=vlmblue, right=0.7cm of consensus] (dimgen) {VLM Dimension\\Generator};
\node[data, right=0.7cm of dimgen, text width=3.6cm] (dims) {\footnotesize $\mathcal{D}_c$: \textit{Fa\c{c}ade Quality},\\[-2pt]\textit{Vegetation}, \textit{Pavement},\\[-2pt]\textit{Vehicle}, \textit{Modernity},\\[-2pt]\textit{Infrastructure}, \textit{Cleanliness}};

\draw[arr=vlmblue] (consensus) -- (dimgen);
\draw[arr=vlmblue] (dimgen) -- (dims);

% ========== STAGE 2 ==========
\node[stage=vlmblue, text width=3.2cm, below=1.4cm of s1label] (s2label) {\color{vlmblue}Stage 2\\[-1pt]\footnotesize Multi-Agent Feature\\[-2pt]\footnotesize Distillation};

% Input images
\node[data, right=0.6cm of s2label, text width=1.6cm] (imgpair) {Image Pair\\$(x_A, x_B)$};

% Multi-agent chain
\node[agent, right=0.7cm of imgpair] (obs) {Observer\\{\tiny $\tau{=}0.3$}};
\node[agent, right=0.4cm of obs] (deb) {Debater\\{\tiny $\tau{=}0.5$}};
\node[agent, right=0.4cm of deb] (jdg) {Judge\\{\tiny $\tau{=}0.1$}};

\draw[arr=vlmblue] (imgpair) -- (obs);
\draw[arr=vlmblue] (obs) -- (deb);
\draw[arr=vlmblue] (deb) -- (jdg);

% Semantic scores output
\node[data, right=0.5cm of jdg, text width=1.8cm] (semscore) {$S(x)\in[1,10]^7$\\{\tiny Semantic Scores}};
\draw[arr=vlmblue] (jdg) -- (semscore);

% CLIP stream
\node[block=clipteal, below=0.6cm of imgpair, text width=1.8cm] (clip) {Frozen CLIP\\{\footnotesize ViT-L/14}};
\node[data, right=2.3cm of clip, text width=2.0cm] (clipfeat) {$\phi(x)\in\mathbb{R}^{768}$\\{\tiny Visual Embed.}};
\draw[arr=clipteal] (imgpair.south) |- (clip);
\draw[arr=clipteal] (clip) -- (clipfeat);

% Fusion
\node[block=vlmblue, right=0.6cm of clipfeat, text width=3.2cm] (fusion) {Hybrid Vector Fusion\\$h{=}[\alpha\cdot\bar\phi,\;(1{-}\alpha)\cdot\bar{S}]$};

\draw[arr=vlmblue] (semscore.south) |- ([yshift=-0.1cm]fusion.north) -- (fusion.north);
\draw[arr=clipteal] (clipfeat) -- (fusion);

% Dimensions feed down
\draw[darr=vlmblue] (dims.south) -- ++(0,-0.5cm) -| ([xshift=-0.3cm]obs.north);

% ========== STAGE 3 ==========
\node[stage=vlmblue, text width=3.2cm, below=2.6cm of s2label] (s3label) {\color{vlmblue}Stage 3\\[-1pt]\footnotesize Local Manifold\\[-2pt]\footnotesize Calibration (LWRR)};

% Ref/Pool split
\node[data, right=0.6cm of s3label, text width=1.5cm] (refsplit) {$\mathcal{D}_{\mathrm{ref}}$\\{\tiny Human Labels}\\{\tiny + TrueSkill}};
\node[data, below=0.3cm of refsplit, text width=1.5cm] (poolsplit) {$\mathcal{D}_{\mathrm{pool}}$\\{\tiny To Calibrate}};

% Hybrid space
\node[block=vlmblue, right=0.8cm of refsplit, text width=2.2cm, yshift=-0.4cm] (hybdiff) {$\Delta_{\mathrm{hybrid}}(A,B)$\\{\tiny Differential Space}};

\draw[arr=gray] (refsplit) -- (hybdiff);
\draw[arr=gray] (poolsplit) -- (hybdiff);

% KNN
\node[block=clipteal, right=0.6cm of hybdiff, text width=1.6cm] (knn) {$K$-NN Search\\{\tiny $K{=}20$}};
\draw[arr=vlmblue] (hybdiff) -- (knn);

% LWRR
\node[block=vlmblue, right=0.6cm of knn, text width=2.6cm] (lwrr) {Weighted Ridge\\Regression\\{\tiny $\hat{\mathbf{w}}{=}(X^\top WX{+}\lambda I)^{-1}X^\top Wy$}};
\draw[arr=vlmblue] (knn) -- (lwrr);

% Output
\node[data, right=0.6cm of lwrr, text width=1.8cm] (output) {$\hat{y}$: Calibrated\\Prediction\\{\tiny + $R^2$ Audit}};
\draw[arr=vlmblue] (lwrr) -- (output);

% Connect fusion to stage 3
\draw[arr=vlmblue] (fusion.south) -- ++(0,-0.5cm) -| (hybdiff.north);

% Data isolation annotation
\node[below=0.15cm of poolsplit, font=\tiny\itshape, text=feedred] {$\mathcal{D}_{\mathrm{ref}} \cap \mathcal{D}_{\mathrm{pool}} = \varnothing$};

% Stage labels on left
\begin{scope}[on background layer]
  \node[fit=(consensus)(dims), draw=vlmblue!30, fill=vlmblue!3, rounded corners=5pt, inner sep=4pt] {};
  \node[fit=(imgpair)(semscore)(clip)(fusion), draw=vlmblue!30, fill=vlmblue!3, rounded corners=5pt, inner sep=4pt] {};
  \node[fit=(refsplit)(poolsplit)(output), draw=vlmblue!30, fill=vlmblue!3, rounded corners=5pt, inner sep=4pt] {};
\end{scope}

\end{tikzpicture}%
}
\caption{Overview of the UrbanAlign framework.
\textbf{Stage~1} discovers interpretable semantic dimensions from high/low-consensus exemplars via a VLM.
\textbf{Stage~2} distils VLM knowledge into continuous dimension scores via the Observer--Debater--Judge chain and fuses them with CLIP embeddings into hybrid vectors ($\alpha{=}0.3$).
\textbf{Stage~3} calibrates predictions by locally-weighted ridge regression (LWRR) on the hybrid differential manifold, with $R^2$ as an interpretability audit.
Data isolation: $\mathcal{D}_{\mathrm{ref}} \cap \mathcal{D}_{\mathrm{pool}} = \varnothing$.}
\label{fig:framework}
\end{figure*}

\subsection{Problem Formulation}
\label{sec:formulation}

Given an urban image set $\mathcal{X}{=}\{x_1,\ldots,x_N\}$ and a perception category $d$ (\eg, wealthy), we aim to predict human pairwise preferences:
\begin{equation}
y_{ij} =
\begin{cases}
\text{left}  & \text{if } x_i \succ x_j \text{ on dimension } d,\\
\text{right} & \text{if } x_i \prec x_j,\\
\text{equal} & \text{if } x_i \approx x_j.
\end{cases}
\label{eq:pairwise}
\end{equation}
The labelled pool $\mathcal{D}_{L}{=}\{(x_i,x_j,y_{ij})\}$ contains crowdsourced pairwise comparisons (from Place Pulse~2.0).
We split $\mathcal{D}_{L}$ into a \emph{reference set} $\mathcal{D}_{\mathrm{ref}}$ (with human labels and TrueSkill ratings) and a \emph{pool set} $\mathcal{D}_{\mathrm{pool}}$ (to be calibrated), with strict isolation: $\mathcal{D}_{\mathrm{ref}} \cap \mathcal{D}_{\mathrm{pool}} = \varnothing$.

\noindent\textbf{Notation.}
$\phi{:}\,\mathcal{X}{\to}\mathbb{R}^{768}$ denotes a pre-trained CLIP encoder (ViT-L/14).
TrueSkill~\cite{herbrich2006trueskill} converts pairwise comparisons to continuous ratings $\mu_i$ per image.

% ---------------------------------------------------------------
\subsection{Stage 1: Semantic Dimension Extraction}
\label{sec:stage1}

\noindent\textbf{Motivation.}
Instead of asking a VLM ``Which image looks more wealthy?'' (a poorly calibrated end-to-end query), we first decompose the abstract percept into interpretable, continuously scorable sub-dimensions.
This is analogous to defining the concept set in a CBM~\cite{koh2020concept}: the dimensions form an interpretable bottleneck through which all subsequent predictions must pass.

\smallskip\noindent\textbf{TrueSkill consensus sampling.}
We convert pairwise comparisons to continuous scores via TrueSkill ($\mu{=}25$, $\sigma{=}8.33$, draw probability 0.10).
We sample \emph{high-consensus} ($\mu{>}\tau_{h},\;\sigma{<}\sigma_{\mathrm{med}}$) and \emph{low-consensus} ($\mu{<}\tau_{l},\;\sigma{<}\sigma_{\mathrm{med}}$) exemplars, where $\tau_{h}$ and $\tau_{l}$ are the 75th and 25th percentiles.
Five images per group form the consensus set $\mathcal{S}_{\mathrm{consensus}}$ (10 total).

\smallskip\noindent\textbf{Dimension extraction.}
A structured prompt presents the high/low exemplars with their TrueSkill scores and PCA-compressed CLIP embeddings (8D).
The VLM outputs $n\in[5,10]$ dimensions in JSON, each containing a name, description, and high/low indicators.
For the \emph{wealthy} category, the extracted dimensions are:
\emph{Fa\c{c}ade Quality}, \emph{Vegetation Maintenance}, \emph{Pavement Integrity}, \emph{Vehicle Quality}, \emph{Building Modernity}, \emph{Infrastructure Condition}, and \emph{Street Cleanliness}.

These dimensions satisfy three requirements: \emph{visual observability} (scorable from street-view images), \emph{measurability} (1--10 continuous scale), and \emph{universality} (applicable across cities).

% ---------------------------------------------------------------
\subsection{Stage 2: Multi-Agent Feature Distillation}
\label{sec:stage2}

\noindent\textbf{From classification to feature extraction.}
Rather than querying a VLM for a one-shot discrete label, we extract a \emph{hybrid feature vector}:
\begin{equation}
h(x) = [\phi_{\text{CLIP}}(x),\; S(x)],
\label{eq:hybrid}
\end{equation}
where $S(x)\in[1,10]^{n}$ is the vector of semantic dimension scores produced by the VLM.
This hybrid representation fuses low-level visual patterns (CLIP) with high-level domain-specific semantics---analogous to augmenting a neural backbone with concept scores in a post-hoc CBM~\cite{yuksekgonul2023posthoc}.

\smallskip\noindent\textbf{$2{\times}2$ factorial design.}
To systematically evaluate component contributions, we define four synthesis modes:
\begin{center}
\small
\begin{tabular}{l|cc}
\toprule
 & Single-Shot & Multi-Agent \\
\midrule
Single-Image & Mode~1 & Mode~3 \\
Pairwise     & Mode~2 & Mode~4 \\
\bottomrule
\end{tabular}
\end{center}
\emph{Factor 1 (Pairwise context):} Modes~2/4 present both images simultaneously, enabling relative comparison.
\emph{Factor 2 (Multi-agent reasoning):} Modes~3/4 employ the deliberation chain below.

\smallskip\noindent\textbf{Observer--Debater--Judge deliberation.}
Inspired by multi-agent debate~\cite{du2024improving,liang2024encouraging} and self-consistency~\cite{wang2023selfconsistency}, we decompose reasoning into three agents (all powered by GPT-4o with different temperatures):

\begin{itemize}
\item \textbf{Observer} ($\tau{=}0.3$): describes visually observable details per dimension, without judgement---suppressing confirmation bias.
\item \textbf{Debater} ($\tau{=}0.5$): argues both for HIGH and LOW scores on each dimension---exploring opposing perspectives.
\item \textbf{Judge} ($\tau{=}0.1$): synthesises descriptions and arguments to produce final scores $S(x)\in[1,10]^{n}$.
\end{itemize}

\smallskip\noindent\textbf{Intensity-based winner determination.}
For each pair $(x_A, x_B)$, we also record the \emph{intensity}: $I{=}|S(x_A) - S(x_B)|$ summed across dimensions.
Pairs with intensity below a threshold $\sigma_I$ are labelled ``equal''; otherwise the VLM's predicted winner is retained.

% ---------------------------------------------------------------
\subsection{Stage 3: Local Manifold Calibration (LWRR)}
\label{sec:stage3}

This stage is the core algorithmic contribution: a post-hoc calibration layer analogous to the linear predictor in a CBM~\cite{koh2020concept}, but \emph{locally adaptive} to account for heterogeneous perception patterns across the visual-semantic manifold.

\subsubsection{Hybrid Differential Space.}

CLIP, pre-trained on web-scale image-text pairs, provides strong general features but exhibits distribution shift on street-view scenes.
We augment it with semantic dimension scores to form the hybrid differential encoding:
\begin{equation}
\Delta_{\text{hybrid}}(A,B)=\bigl[\alpha\cdot\overline{\Delta}_{\text{CLIP}}(A,B),\;(1{-}\alpha)\cdot\overline{\Delta}_{\text{Sem}}(A,B)\bigr],
\label{eq:hybrid_diff}
\end{equation}
where $\overline{\Delta}_{\text{CLIP}}$ uses L2-normalised CLIP embeddings, $\Delta_{\text{Sem}}{=}S(A){-}S(B)$ is normalised to $[0,1]$ by dividing by 10, and $\alpha\in[0,1]$ balances visual and semantic components.
We find $\alpha{=}0.3$ optimal (70\% semantic weight), consistent with the expectation that domain-specific VLM concepts carry more perception-relevant signal than generic CLIP features.

\subsubsection{LWRR Algorithm.}

\noindent\textbf{Step~1: Build reference manifold.}
From $\mathcal{D}_{\mathrm{ref}}$, compute hybrid differential vectors and TrueSkill score differences $y_i^{\text{TS}}{=}\mu_A{-}\mu_B$.
Mirror augmentation doubles the reference set: for each $(A,B,y)$, add $(B,A,-y)$.

\noindent\textbf{Step~2: Neighbourhood search.}
For query $\Delta_q^{\mathrm{hybrid}} \in \mathcal{D}_{\mathrm{pool}}$, compute cosine similarities to all reference points and select the $K$ nearest neighbours $\mathcal{N}_K$ ($K{=}20$).

\noindent\textbf{Step~3: Kernel weighting.}
$w_i=\exp(s_i/\tau)$ for $i\in\mathcal{N}_K$, where $s_i$ is cosine similarity and $\tau{=}1.0$ is kernel bandwidth.

\noindent\textbf{Step~4: Local ridge regression.}
Solve for local dimension weights:
\begin{equation}
\hat{\mathbf{w}} = \arg\min_{\mathbf{w}} \sum_{i\in\mathcal{N}_K} w_i\bigl(y_i^{\text{TS}}-\mathbf{w}^{\!\top}\Delta_{\text{Sem},i}\bigr)^{2} + \lambda\|\mathbf{w}\|^{2},
\label{eq:lwrr}
\end{equation}
where $\lambda{=}1.0$ is the ridge penalty.
The solution is $\hat{\mathbf{w}}=(X^{\!\top}WX+\lambda I)^{-1}X^{\!\top}Wy$, where $W{=}\mathrm{diag}(w_1,\ldots,w_K)$ and $X$ stacks the semantic differentials.

\noindent\textbf{Step~5: Prediction and audit.}
\begin{equation}
\hat{\delta}_q = \hat{\mathbf{w}}^{\!\top}\Delta_{\text{Sem}}^{q},\qquad
R^{2} = 1-\frac{\sum_{i}w_i(y_i^{\text{TS}}-\hat{y}_i)^{2}}{\sum_{i}w_i(y_i^{\text{TS}}-\bar{y})^{2}}.
\label{eq:r2}
\end{equation}
When $R^{2}$ is high, abstract perception in this neighbourhood is linearly explainable by the mid-level dimensions.

\noindent\textbf{Step~6: Statistical re-inference.}
\begin{equation}
\hat{y} =
\begin{cases}
\text{equal} & \text{if } |\hat{\delta}|<\varepsilon \text{ or } c_{\mathrm{eq}}>\theta, \\
\text{left}  & \text{if } \hat{\delta}>0, \\
\text{right} & \text{otherwise},
\end{cases}
\label{eq:reinfer}
\end{equation}
where $\varepsilon{=}0.8$ is the score-difference threshold and $c_{\mathrm{eq}}{=}|\{i\in\mathcal{N}_K: l_i{=}\text{equal}\}|/K$ measures local equal-consensus.

\subsubsection{Why Post-hoc Calibration Works.}

The theoretical motivation is straightforward.
VLMs are powerful feature extractors: their dimension scores $S(x)$ capture genuine perceptual signals (we show up to 64.4\% discriminative power per dimension).
However, the VLM's \emph{own} mapping from $S$ to discrete labels is poorly calibrated---it lacks exposure to this specific population's decision boundaries.
LWRR learns these boundaries locally from a small reference set.
The procedure is analogous to a post-hoc CBM~\cite{yuksekgonul2023posthoc}: concept activations (our dimension scores) are extracted from a frozen backbone (our VLM), then a lightweight predictor (our LWRR) is trained on labelled data.
The key advantage of LWRR over a single global linear layer is \emph{local adaptivity}: different regions of the perception manifold may weight dimensions differently (\eg, ``vegetation maintenance'' may matter more in suburban neighbourhoods than in downtown areas).

\subsection{Stage~4: End-to-End Dimension Optimization}
\label{sec:e2e_method}

The semantic dimensions discovered in Stage~1 depend on the VLM's sampling temperature and the particular consensus exemplars presented.
We close the loop with an automated search over dimension sets.

\smallskip\noindent\textbf{Formulation.}
Let $\mathcal{D}_c^{(t)}$ denote the dimension set generated at trial $t$ with temperature $\tau_{\mathrm{gen}}^{(t)} = 0.7 + 0.1t$.
Each trial runs the full pipeline---dimension extraction $\to$ pairwise scoring (Mode~2) $\to$ LWRR calibration---on a fixed evaluation sample.
The optimal set is:
\begin{equation}
\mathcal{D}_c^{*} = \arg\max_{t \in \{1,\ldots,T\}} \;\mathrm{Acc}\bigl(\text{LWRR}\bigl(\text{Score}(\mathcal{D}_{\mathrm{pool}},\, \mathcal{D}_c^{(t)}),\, \mathcal{D}_{\mathrm{ref}}\bigr)\bigr).
\label{eq:e2e}
\end{equation}

\smallskip\noindent\textbf{Elite seeding.}
Trials after the first receive the current best $\mathcal{D}_c^{*}$ as a soft reference in the generation prompt, encouraging the VLM to produce informed variants rather than independent samples.
This converts the search from uniform sampling into guided refinement.

\smallskip\noindent\textbf{Convergence.}
If each trial draws independently from the VLM's distribution over an effective dimension-set space $\mathcal{F}_{\mathrm{eff}}$, the probability of missing $\mathcal{D}_c^*$ after $T$ trials satisfies:
\begin{equation}
\Pr\bigl[\mathcal{D}_c^* \notin \{\mathcal{D}_c^{(1)}, \ldots, \mathcal{D}_c^{(T)}\}\bigr] \leq (1 - p^*)^T,
\label{eq:convergence}
\end{equation}
where $p^*$ is the single-trial probability of generating the optimal set.
The progressive temperature schedule traverses the VLM's generative distribution from its mode (exploitation) toward higher-entropy regions (exploration), analogous to random search with a guided prior~\cite{bergstra2012random}.
Full theoretical analysis---information-bottleneck interpretation, sample complexity, and connection to LWRR convergence---is provided in supplementary Section~E.

\smallskip\noindent\textbf{Complete algorithm.}
The full pseudocode is provided in supplementary Section~G.
In summary: Stage~1 extracts $n$ semantic dimensions; Stage~2 scores all pairs via Observer--Debater--Judge; Stage~3 calibrates via LWRR; and Stage~4 optimizes the dimension set end-to-end.


% ===============================================================
% 4  EXPERIMENTS
% ===============================================================
\section{Experiments}
\label{sec:exp}

\subsection{Setup}

\noindent\textbf{Dataset.}
Place Pulse 2.0~\cite{salesses2013collaborative}: 110,688 street-view images from 56 cities with 1,169,078 pairwise comparisons across six perception categories (safety, lively, beautiful, wealthy, depressing, boring).
We report results across \textbf{all six categories}, with detailed ablation on \emph{wealthy}.
Per category, we sample ${\sim}$256--288 pairs; after Stage~3 alignment on the pool split, 65--87 pairs remain (after excluding human ``equal'' labels) for evaluation.

\smallskip\noindent\textbf{Data protocol.}
$\mathcal{D}_{\mathrm{ref}}$ (70\% of sampled pairs) provides human labels and TrueSkill ratings for LWRR calibration; $\mathcal{D}_{\mathrm{pool}}$ (30\%) is calibrated and evaluated.
$\mathcal{D}_{\mathrm{ref}} \cap \mathcal{D}_{\mathrm{pool}} = \varnothing$ (random split, seed\,=\,42).

\smallskip\noindent\textbf{Implementation.}
CLIP ViT-L/14 (768-d) for visual features.
GPT-4o (\texttt{gpt-4o}) for all VLM calls.
Stage~2: Observer $\tau{=}0.3$, Debater $\tau{=}0.5$, Judge $\tau{=}0.1$; single-shot modes $\tau{=}0.0$.
Stage~3: $K{=}20$, $\tau_{\mathrm{kernel}}{=}1.0$, $\lambda{=}1.0$, $\varepsilon{=}0.8$, $\theta{=}0.6$, $\alpha{=}0.3$.

\smallskip\noindent\textbf{Baselines.}
(C0)~End-to-end Siamese network on ResNet-50 features~\cite{koch2015siamese};
(C1)~End-to-end Siamese network on CLIP features;
(C2)~Segmentation-feature regression via DeepLabV3~\cite{chen2018deeplabv3plus} (21-class semantic segmentation $\to$ 8 urban-group pixel proportions $\to$ logistic regression);
(C3)~Zero-shot VLM: GPT-4o with a simple prompt ``Which looks more \texttt{<category>}?'' without semantic dimensions or multi-agent reasoning.

\smallskip\noindent\textbf{Metrics.}
Accuracy (Acc), Cohen's $\kappa$ (chance-corrected agreement), F1 (macro).
We report metrics both including and excluding pairs where the human label is ``equal'' (as these are inherently ambiguous).

% ---------------------------------------------------------------
\subsection{Main Results}

\begin{table}[t]
\centering
\caption{Comparison on Place Pulse 2.0 across all six perception categories. Accuracy and $\kappa$ are reported excluding ``equal'' human labels. UrbanAlign (Mode~4 + LWRR) outperforms all baselines on average (61.3\% vs.\ 59.1\% for the best baseline).}
\label{tab:main}
\small
\begin{tabular}{l c c c c c c}
\toprule
 & \multicolumn{2}{c}{Safety} & \multicolumn{2}{c}{Beautiful} & \multicolumn{2}{c}{Lively} \\
\cmidrule(lr){2-3} \cmidrule(lr){4-5} \cmidrule(lr){6-7}
Method & Acc.(\%) & $\kappa$ & Acc.(\%) & $\kappa$ & Acc.(\%) & $\kappa$ \\
\midrule
(C0) ResNet50 Siamese & 61.9 & 0.20 & 57.1 & 0.14 & 53.1 & 0.06 \\
(C1) CLIP Siamese & 59.5 & 0.17 & \textbf{62.9} & \textbf{0.26} & 49.4 & $-$0.03 \\
(C2) Seg.\ Regression & \textbf{64.3} & \textbf{0.28} & \textbf{62.9} & 0.25 & 54.3 & 0.10 \\
(C3) GPT-4o zero-shot & 31.0 & 0.14 & 34.3 & 0.09 & 46.9 & $-$0.03 \\
UrbanAlign (Mode~4 + LWRR) & \textbf{66.3} & \textbf{0.30} & 56.5 & 0.13 & \textbf{57.1} & \textbf{0.15} \\
\midrule
 & \multicolumn{2}{c}{Wealthy} & \multicolumn{2}{c}{Boring} & \multicolumn{2}{c}{Depressing} \\
\cmidrule(lr){2-3} \cmidrule(lr){4-5} \cmidrule(lr){6-7}
Method & Acc.(\%) & $\kappa$ & Acc.(\%) & $\kappa$ & Acc.(\%) & $\kappa$ \\
\midrule
(C0) ResNet50 Siamese & 53.0 & 0.09 & 36.5 & $-$0.20 & \textbf{62.7} & \textbf{0.26} \\
(C1) CLIP Siamese & 62.7 & 0.26 & 47.3 & $-$0.05 & 53.3 & 0.07 \\
(C2) Seg.\ Regression & 57.8 & 0.17 & \textbf{60.8} & \textbf{0.25} & 54.7 & 0.09 \\
(C3) GPT-4o zero-shot & 63.9 & 0.33 & 43.2 & 0.06 & 41.3 & 0.03 \\
UrbanAlign (Mode~4 + LWRR) & \textbf{72.3} & \textbf{0.44} & 61.5 & 0.24 & 54.3 & 0.09 \\
\bottomrule
\end{tabular}
\\[4pt]
\footnotesize{\textsuperscript{*}Aligned modes evaluated on disjoint pool subset (65--83 pairs after excl.\ equal per category). Baselines evaluated on the same pool split.}
\end{table}

\Cref{tab:main} and \cref{fig:comparison} present the main comparison.
Key findings:

(1)~UrbanAlign (Mode~4 + LWRR) achieves 61.3\% accuracy ($\kappa{=}0.22$) on average across all six categories, with 72.3\% accuracy ($\kappa{=}0.44$) on \emph{wealthy} perception being the strongest.
This represents a \textbf{+17.9\,pp} average gain over the zero-shot VLM baseline (C3: 43.4\%) and outperforms the best traditional baseline (C2 Segmentation Regression: 59.1\%) by +2.2\,pp.

(2)~The zero-shot VLM (C3) shows large variance across categories (31.0\%--63.9\% accuracy), confirming VLMs possess strong perceptual capabilities for some perceptions (\eg, \emph{wealthy}: 63.9\%) but struggle with others (\eg, \emph{safety}: 31.0\%).

(3)~LWRR calibration provides a +7.0\,pp average boost over Mode~4 Raw (54.3\%$\to$61.3\%).

% ---------------------------------------------------------------
\subsection{Ablation: $2{\times}2$ Factorial Design}

\begin{table}[t]
\centering
\caption{$2{\times}2$ factorial ablation on \emph{wealthy} (accuracy \%\,/\,$\kappa$, post-LWRR alignment). Multi-agent reasoning is effective only when combined with pairwise context (Mode~4); alone (Mode~3) it hurts performance.}
\label{tab:factorial}
\small
\begin{tabular}{l|cc|c}
\toprule
 & Single-Shot & Multi-Agent & Multi-Agent Gain \\
\midrule
Single-Image & 59.0\,/\,0.18 & 54.1\,/\,0.08 & $-$4.9 \\
Pairwise     & 64.2\,/\,0.28 & \textbf{72.3}\,/\,\textbf{0.44} & +8.1 \\
\midrule
Pairwise Gain & +5.2 & +18.2 & --- \\
\bottomrule
\end{tabular}
\end{table}

\begin{table}[t]
\centering
\caption{LWRR calibration gain (Mode~4, accuracy \% excl.\ ``equal'') across all six categories. LWRR improves four of six categories; the largest gain is on \emph{boring} (+23.3\,pp).}
\label{tab:vrm_gain_main}
\small
\begin{tabular}{lcccccc|c}
\toprule
 & Safety & Beaut. & Lively & Wealthy & Boring & Depress. & Avg \\
\midrule
Raw Acc.\ (\%)     & 59.6 & 60.2 & 59.1 & 65.0 & 38.2 & 43.6 & 54.3 \\
Aligned Acc.\ (\%) & 66.3 & 56.5 & 57.1 & 72.3 & 61.5 & 54.3 & 61.3 \\
$\Delta$ (pp)      & \textbf{+6.7} & $-$3.6 & $-$2.0 & \textbf{+7.3} & \textbf{+23.3} & \textbf{+10.7} & \textbf{+7.0} \\
\midrule
Raw $\kappa$       & 0.19 & 0.21 & 0.18 & 0.30 & $-$0.23 & $-$0.13 & 0.09 \\
Aligned $\kappa$   & 0.30 & 0.13 & 0.15 & 0.44 & 0.24 & 0.09 & 0.22 \\
\midrule
$n$ (pool pairs)   & 83 & 69 & 77 & 83 & 65 & 70 & --- \\
\bottomrule
\end{tabular}
\end{table}

\Cref{tab:factorial} shows the post-alignment $2{\times}2$ factorial results:

\smallskip\noindent\textbf{Pairwise context (Modes~2/4 vs.~1/3).}
Across single-shot, pairwise input improves accuracy by +5.2\,pp (Mode~1$\to$2).
Across multi-agent, the gain is much larger: +18.2\,pp (Mode~3$\to$4).

\smallskip\noindent\textbf{Multi-agent reasoning (Modes~3/4 vs.~1/2).}
Multi-agent reasoning \emph{hurts} in the single-image setting ($-$4.9\,pp, Mode~1$\to$3) but provides a large gain in the pairwise setting (+8.1\,pp, Mode~2$\to$4).
This is a key finding: the Observer--Debater--Judge chain requires pairwise context to function effectively.
With single images, the Debater's ``argue both sides'' strategy leads to increased noise rather than constructive deliberation; with paired images providing concrete comparative anchors, the deliberation chain generates genuinely complementary perspectives.

\smallskip\noindent\textbf{Synergy.}
The two factors interact synergistically rather than additively: Mode~4 gains +13.3\,pp over Mode~1, far more than the sum of individual effects (+5.2 pairwise $+$ ($-$4.9) multi-agent $=$ +0.3).
This synergy arises because multi-agent debate is \emph{enabled by} pairwise comparison---arguing about relative merits requires both items to be present.

\smallskip\noindent\textbf{LWRR alignment gain.}
\Cref{tab:vrm_gain_main} shows LWRR calibration gain across all six categories for Mode~4.
LWRR yields a +7.0\,pp average improvement (54.3\%$\to$61.3\%), with the most dramatic effect on \emph{boring} (+23.3\,pp) and \emph{depressing} (+10.7\,pp)---categories where the VLM's raw predictions perform near chance (38.2\% and 43.6\%).
On \emph{wealthy}, only Mode~4 benefits from LWRR calibration (+7.3\,pp); other modes show zero or negative gain (Mode~1: $-$10.1\,pp, Mode~2: $\pm$0.0\,pp, Mode~3: $-$11.7\,pp), confirming that self-consistent multi-agent pairwise scores are prerequisite for meaningful local calibration (full per-mode breakdown in supplementary Section~F).

% ---------------------------------------------------------------
\subsection{Dimension-Level Discriminability}

\begin{table}[t]
\centering
\caption{Per-dimension discriminative power across all six categories (Mode~4, top-3 shown per category). Disc.\ power is the fraction of pool pairs where the dimension score difference correctly predicts the human winner (50\% $=$ chance).}
\label{tab:dim_disc}
\small
\begin{tabular}{llc|llc}
\toprule
\multicolumn{3}{c|}{\textbf{Safety} ($n{=}87$, avg 60.8\%)} & \multicolumn{3}{c}{\textbf{Beautiful} ($n{=}74$, avg 46.9\%)} \\
Dimension & \% && Dimension & \% & \\
\midrule
Landscaping \& Greenery & 70.1 && Natural Elements & 56.8 & \\
Visibility Clarity & 65.5 && Landscape Design & 54.1 & \\
Traffic Conditions & 65.5 && Traffic Organisation & 54.1 & \\
\midrule
\multicolumn{3}{c|}{\textbf{Lively} ($n{=}87$, avg 51.5\%)} & \multicolumn{3}{c}{\textbf{Wealthy} ($n{=}87$, avg 60.4\%)} \\
Dimension & \% && Dimension & \% & \\
\midrule
Temporal Dynamism & 55.2 && Fa\c{c}ade Quality & 64.4 & \\
Human Activity & 54.0 && Vegetation Maint. & 62.1 & \\
Street Furniture & 54.0 && Pavement Integrity & 60.9 & \\
\midrule
\multicolumn{3}{c|}{\textbf{Boring} ($n{=}80$, avg 34.5\%)} & \multicolumn{3}{c}{\textbf{Depressing} ($n{=}78$, avg 47.1\%)} \\
Dimension & \% && Dimension & \% & \\
\midrule
Visual Monotony & 43.8 && Social Activity & 53.8 & \\
Activity Sparsity & 41.2 && Arch.\ Diversity & 53.8 & \\
Functional Diversity & 38.8 && Building Utilisation & 51.3 & \\
\bottomrule
\end{tabular}
\\[4pt]
\footnotesize{Full per-dimension tables for all categories are provided in the supplementary material.}
\end{table}

\Cref{tab:dim_disc} reports per-dimension discriminative power for Mode~4 across all six categories.
Performance varies substantially: \emph{safety} (avg.\ 60.8\%) and \emph{wealthy} (avg.\ 60.4\%) achieve the highest discriminability, indicating that the VLM discovers dimensions well-aligned with human perception for these categories.
\emph{Boring} (avg.\ 34.5\%) is the weakest, with most dimensions below chance, suggesting that ``boringness'' is harder to decompose into visually scorable sub-dimensions.

For \emph{wealthy}, \emph{Fa\c{c}ade Quality} (64.4\%) and \emph{Vegetation Maintenance} (62.1\%) are the strongest predictors, consistent with urban planning theory that building maintenance and landscaping signal neighbourhood affluence~\cite{wilson1982broken}.
The aggregate combination via LWRR (72.3\%) substantially exceeds any single dimension, validating the multi-dimensional bottleneck.

\smallskip\noindent\textbf{Interpretability example.}
For a pair scored by LWRR with high $R^{2}$, the local weights $\hat{\mathbf{w}}$ reveal: Fa\c{c}ade Quality ($w{=}0.28$), Pavement Integrity ($w{=}0.23$), Vehicle Quality ($w{=}0.19$) as the top contributors.
This is directly actionable: planners can target fa\c{c}ade renovation and pavement repair in areas perceived as least wealthy.

% ---------------------------------------------------------------
\subsection{Sensitivity Analysis}

We study the sensitivity of Stage~3 to key hyperparameters (\cref{tab:sensitivity}).
All experiments vary one parameter while holding others at default.

\begin{table}[t]
\centering
\caption{Sensitivity of Mode~4 alignment accuracy (\%, excluding ``equal'') to key Stage~3 hyperparameters, per category. Bold indicates the best value for each category.}
\label{tab:sensitivity}
\small
\setlength{\tabcolsep}{3.5pt}
\begin{tabular}{lcccccc|c}
\toprule
& Safety & Beaut. & Lively & Wealthy & Boring & Depress. & Avg \\
\midrule
\multicolumn{8}{l}{\textit{$K_{\max}$ (neighbourhood size)}} \\
\quad 10 & 64.6 & \textbf{59.1} & 45.0 & 57.5 & 52.9 & 42.0 & 53.5 \\
\quad 20 (default) & 66.3 & 56.5 & \textbf{57.1} & \textbf{72.3} & \textbf{61.5} & 54.3 & \textbf{61.3} \\
\quad 30 & 63.7 & 52.2 & 56.9 & 66.7 & 59.1 & \textbf{55.1} & 58.9 \\
\quad 50 & \textbf{67.5} & 56.2 & 53.4 & 67.9 & 51.6 & 54.0 & 58.4 \\
\midrule
\multicolumn{8}{l}{\textit{$\alpha$ (CLIP vs.\ semantic weight)}} \\
\quad 0.1 & 65.4 & 57.6 & 51.3 & 67.5 & 57.4 & 55.1 & 59.1 \\
\quad 0.3 (default) & \textbf{66.3} & 56.5 & \textbf{57.1} & \textbf{72.3} & \textbf{61.5} & 54.3 & \textbf{61.3} \\
\quad 0.5 & 65.4 & \textbf{63.3} & 52.7 & 65.8 & 54.3 & \textbf{60.0} & 60.3 \\
\quad 0.7 & 64.1 & 59.1 & 56.8 & 64.2 & 50.7 & 45.6 & 56.7 \\
\midrule
\multicolumn{8}{l}{\textit{SELECTION\_RATIO (confidence filter)}} \\
\quad 0.6 & \textbf{80.8} & \textbf{68.3} & \textbf{63.3} & \textbf{75.0} & \textbf{65.5} & \textbf{57.1} & \textbf{68.3} \\
\quad 0.8 & 71.6 & 58.2 & 58.7 & 74.6 & 65.5 & 56.1 & 64.1 \\
\quad 0.9 & 68.4 & 58.1 & 58.3 & 72.4 & 61.5 & 53.8 & 62.1 \\
\quad 1.0 (default) & 66.3 & 56.5 & 57.1 & 72.3 & 61.5 & 54.3 & 61.3 \\
\bottomrule
\end{tabular}
\\[4pt]
\footnotesize{$\tau_{\mathrm{kernel}}$ and $\lambda_{\mathrm{ridge}}$ are omitted ($<$0.5\,pp avg.\ variation). Full sweeps including all parameters in supplementary Section~C.}
\end{table}

\noindent\textbf{SELECTION\_RATIO} is the most impactful parameter: reducing from 1.0 to 0.6 yields a +7.0\,pp average gain (61.3\%$\to$68.3\%), with \emph{safety} showing the largest individual improvement (+14.5\,pp to 80.8\%).
The effect is consistent across all six categories, confirming that confidence-based filtering of aligned pairs substantially improves output quality.

\noindent\textbf{$K_{\max}$} is the second most sensitive: $K{=}10$ collapses the average to 53.5\% while $K{=}20$ achieves 61.3\%.
Per-category variation reveals that the optimal $K$ is relatively consistent ($K{=}20$ is best for four of six categories), though \emph{safety} favours $K{=}50$ (67.5\%) and \emph{beautiful} favours $K{=}10$ (59.1\%).

\noindent\textbf{$\alpha$ (CLIP vs.\ semantic weight)} shows a clear average optimum at $\alpha{=}0.3$ (70\% semantic weight), though per-category optima differ: \emph{beautiful} and \emph{depressing} favour $\alpha{=}0.5$ (63.3\% and 60.0\% respectively), suggesting that CLIP features carry more useful information for aesthetics-related categories.
This parallels the CBM finding that concept-based predictions can match or exceed end-to-end models when concepts are well-chosen~\cite{koh2020concept}.

\noindent\textbf{$\tau_{\mathrm{kernel}}$ and $\lambda_{\mathrm{ridge}}$} have negligible effect ($<$0.5\,pp variation).
This is theoretically expected: for locally-weighted ridge regression with $K$ neighbours in $n$ dimensions, the excess risk scales as $\mathcal{O}(\lambda + n/(K\lambda))$~\cite{caponnetto2007optimal}, which becomes insensitive to $\lambda$ when $K \gg n$ (here $K{=}20$, $n{=}7$).
Full parameter-sweep grids are provided in the supplementary material (Section~C).

% ---------------------------------------------------------------
\smallskip\noindent\textbf{Dimension optimization.}
As described in Section~\ref{sec:e2e_method}, we automate dimension-set selection via the end-to-end optimization loop (Eq.~\eqref{eq:e2e}).
With $T{=}6$ trials on \emph{wealthy} (temperatures $\tau_{\mathrm{gen}}{=}0.7$--$1.2$), the best dimension set emerges at the lowest temperature ($\tau{=}0.7$, 56.3\% accuracy, $\kappa{=}0.231$), with performance degrading monotonically at higher temperatures.
This suggests that conservative, high-consensus dimensions are most effective, consistent with the $(1{-}p^*)^T$ convergence bound (Eq.~\eqref{eq:convergence}).
Full trial-by-trial results are reported in the supplementary (Section~E).


% ===============================================================
% 5  DISCUSSION
% ===============================================================
\section{Discussion}
\label{sec:discussion}

\noindent\textbf{Why does post-hoc calibration work?}

The central finding is that VLMs are strong \emph{concept extractors} but poor \emph{decision calibrators} for urban perception.
Mode~4 raw dimension scores already contain substantial perceptual information (individual dimensions at 57--64\% discriminative power), yet the VLM's own label predictions are only marginally above the zero-shot baseline.
LWRR solves a simple and intuitive problem: given that dimension scores $S(x)$ are informative, learn the locally-optimal linear mapping from score differences to TrueSkill rating differences.
This is precisely the post-hoc CBM paradigm~\cite{yuksekgonul2023posthoc}: extract concepts from a frozen backbone, then train a lightweight predictor on those concepts.

\smallskip\noindent\textbf{Why local rather than global calibration?}
Urban perception is heterogeneous: the visual cues that signal wealth in a suburban neighbourhood (vegetation, vehicle quality) may differ from those in an urban core (building modernity, infrastructure).
A single global linear model would average over these distinct regimes.
LWRR, by fitting independent weights per neighbourhood, naturally adapts to local manifold geometry.
The sensitivity analysis confirms this: $K_{\max}$ is the most critical parameter, reflecting the tension between local adaptivity and statistical stability.
Theoretically, the LWRR excess risk at optimal regularisation decays as $\mathcal{O}((n/K)^{1/2})$~\cite{caponnetto2007optimal}, explaining why larger $K$ improves accuracy when $K \gg n$; the per-pair weight analysis in supplementary Section~B confirms that dimension importances vary substantially across the manifold, validating the need for local fitting.

\smallskip\noindent\textbf{Connection to the information bottleneck.}
Our pipeline implements a form of information bottleneck~\cite{tishby2000information}: Stage~2 compresses visual input $X$ into dimension scores $S{=}f_{\mathcal{D}_c}(X)\in\mathbb{R}^n$ (a compact bottleneck), and Stage~3 maximises the downstream mutual information $I(\hat{Y}; Y)$ via LWRR.
The $R^{2}$ metric quantifies concept completeness---how much of the TrueSkill variance is explainable by the bottleneck representation.
The end-to-end dimension optimization (Eq.~\eqref{eq:e2e}; supplementary Section~E) formalises this further: selecting the bottleneck that maximises $I(\hat{Y};Y)$ is equivalent to choosing the dimension set with highest alignment accuracy.
Elite seeding biases subsequent trials toward the current-best bottleneck, converting the search from independent sampling into a guided refinement with exponential convergence guarantee.

\smallskip\noindent\textbf{Limitations.}
While we evaluate across all six Place Pulse~2.0 categories, performance varies substantially: \emph{wealthy} (72.3\%) far exceeds \emph{depressing} (54.3\%) and \emph{boring} (61.5\%).
The per-category sample sizes are modest (65--83 pool pairs after filtering).
The current approach assumes static street-view images; temporal and cultural variation remain unexplored.

\smallskip\noindent\textbf{Ethical considerations.}
Automated perception mapping could inform equitable infrastructure investment but also risks reinforcing neighbourhood stereotypes if applied without community engagement.
We recommend use for improvement-oriented planning (identifying streets needing fa\c{c}ade repair, pavement maintenance) rather than discriminatory rankings.


% ===============================================================
% 6  CONCLUSION
% ===============================================================
\section{Conclusion}
\label{sec:conclusion}

We presented \textbf{UrbanAlign}, a post-hoc calibration framework that transforms VLMs from unreliable annotators (43.4\% zero-shot accuracy) into structured perception decoders (61.3\% accuracy on average, $\kappa{=}0.22$) by routing predictions through interpretable semantic dimensions and locally-adaptive calibration.
The framework operationalises the post-hoc concept bottleneck paradigm for VLM-synthesised data: semantic dimensions form a compact information bottleneck, and LWRR provides a locally-adaptive calibration layer whose excess risk is theoretically bounded.

Our $2{\times}2$ factorial ablation reveals that pairwise context and multi-agent reasoning interact synergistically (+13.3\,pp combined gain on \emph{wealthy}), providing guidance for future VLM-based annotation systems.
The per-dimension discriminability analysis confirms that VLM-derived concepts capture genuine perceptual signals (up to 64.4\% per dimension), and the LWRR $R^{2}$ metric provides an interpretability audit that makes each prediction explainable in terms of mid-level semantic dimensions.

Future work will explore cross-city transfer and investigate whether the semantic-dimension--calibration pipeline generalises to other subjective perception tasks such as aesthetic quality assessment, medical image rating, and landscape evaluation.

% ---------------------------------------------------------------
\section*{Acknowledgements}
Omitted for blind review.

% ---- Bibliography ----
\bibliographystyle{splncs04}
\bibliography{main}
\end{document}
